% sec-00: front matter.  Why La Jolla, how the package is organized, the site at
% a glance, the author record and funding position, and how to read the package.
% STAGE 3.  Column widths in the predecessor comparison and the document index
% retuned so no cell wraps to more lines than its neighbours; two sentences
% tightened so section 1 closes on the page it opens.
\docfront{Front Matter}{LA JOLLA MOVE-IN, FRONT MATTER}{Front Matter: The La
Jolla Move-In}

\section{Why La Jolla, and Why Now}

This package stands up one clinical trial site, in La Jolla, San Diego, for one
disease and one protocol. It is the successor to the eleven-document San
Francisco package, which described a general oncology site operating on demand
for twenty-four hours a day and serving 168 unique participants across fifteen
cancer types in a single day with twenty-nine robot instances
\cite{sfsite2026,newtrial2026}.

The La Jolla site is smaller in every dimension that can be counted, and that
narrowing is the point rather than a concession. A Phase 1 trial in resectable
pancreatic ductal adenocarcinoma treats up to eighteen participants across a
3+3 dose escalation \cite{phase1protocol}. A facility sized for 168 participants a
day would spend its federal award on capacity the protocol forbids it to use.
Every requirement in the thirteen documents that follow the legislative part is
derived from the protocol rather than from a service-line planning ratio.

The choice of La Jolla follows from the partner. The intended partner of choice
at the feasibility stage is UC San Diego Moores Cancer Center, which sits in the
same community plan area, roughly 1.6 miles from the leasehold described here.
The twelve-step feasibility sequence, the sponsor-investigator determination, and
the required positioning language are set out in the repository and are carried
into this package without softening. No agreement of any kind exists. UC San
Diego is not described here as a partner, a sponsor, a trial site, or an
endorser, and nothing in this package has been shown to it for approval.

The driver for the move-in is the body of work behind it. Between June 2025 and
August 2026 the author deposited twenty papers that together carry a pancreatic
oncology trial from drug identification through simulation, regulatory
adaptation, protocol, investigational new drug application, and funding
\cite{pdacdeck2026}. Federal policy moved in the same window: the White House
report \emph{Science: A New Golden Age} asks the roughly \$200 billion annual
federal research portfolio to prioritize the individual scientist over legacy
institutions, and its annexed fiscal year 2028 budget priorities memorandum asks
agencies for long-duration grants and non-federal cost share
\cite{kratsios2026sciencegoldenage,kratsiosvought2026fy2028priorities}. This
package is what an independent scientist does with a five-year award once the
policy and the evidence are both in hand.

\begin{mvtable}
\begin{tabularx}{\textwidth}{>{\raggedright\arraybackslash}p{3.3cm} >{\raggedright\arraybackslash}p{4.6cm} >{\raggedright\arraybackslash}X}
\headrow \thead{Dimension} & \thead{San Francisco, 2026} & \thead{La Jolla, this package}\\
\midrule
Scope & All oncology, fifteen cancer types & Pancreatic ductal adenocarcinoma only \\
Operating model & Twenty-four hours, on demand & Extended day, scheduled visits \\
Participants & 168 unique in twenty-four hours & Up to 18 treated, up to 30 screened, over five years \\
Robots & 10 types, 29 instances & 8 types, 14 instances \\
Staff & 8 to 10 full-time equivalents & 11 roles at 3.95 award-funded full-time equivalents \\
Documents & 11 & 15 \\
Funding & Not stated & \$700,000 per year for five years, \$3,500,000 direct \\
Governing evidence & Simulation & Simulation, a filed protocol, and an investigational new drug application \\
\end{tabularx}
\end{mvtable}
\tabcapl{tab:predecessor}{The predecessor and this package, dimension by
dimension. Four rows shrink and three grow, and the growth is in documents,
funding, and evidence rather than in capacity.}

\section{How This Package Is Organized}

Fifteen documents in four parts. Each opens on its own page and can be read
alone. Four documents exist here that had no counterpart in the predecessor,
because the master prompt asks for three things the San Francisco package never
had to answer: conventional pancreatic cancer trial requirements, a lobbying and
federal funding position, and a named staff of eleven.

\begin{mvtable}
\begin{tabularx}{\textwidth}{>{\raggedright\arraybackslash}p{0.7cm} >{\raggedright\arraybackslash}p{1.0cm} >{\raggedright\arraybackslash}p{4.0cm} >{\raggedright\arraybackslash}X}
\headrow \thead{No} & \thead{Part} & \thead{Document} & \thead{What only this document gives the reader}\\
\midrule
01 & I & SB 1188, Site Authorization Act & The state authority under which a site of this kind may exist at all, and the clock the department runs on \\
02 & I & AB 3162, Patient Rights and Robotic Safety Act & The seven participant rights, including the right to a pathway with no model in it \\
03 & I & SB 964, Model Transparency and Data Protection Act & What must be disclosed about the model, and how long a prompt is kept \\
04 & I & H. R. 10412, Federal Act & The federal demonstration designation and the review-timeline reporting requirement \\
05 & II & San Diego Municipal Code Update & Whether this use is permitted on this parcel, and what permit carries it \\
06 & II & Title 22, Division 5, Chapter 15 & The department's application, survey, and enforcement machinery \\
07 & II & FDA LLM and Robotic Workflow Compliance Guide & The requirement-by-requirement map from federal law to a document in this package \\
08 & III & Building Code Standards & What has to be built, from floor loading to the machine room \\
09 & III & Premises Code & Who and what may be where, and the robot envelope numbers \\
10 & III & Parking and Transportation Standards & How a participant arrives, and how one leaves after a Whipple \\
11 & III & Emergency Preparedness Plan & What happens when power, network, model, or building fails \\
12 & IV & Site Activation and Standard Operating Procedures & The five gates and the checklist that opens the site \\
13 & IV & Conventional Trial Requirements Manual & Everything the site needs if the model and the robots were deleted \\
14 & IV & Staffing, Role Delineation, and Move-In Plan & The eleven people, what each is for, and the week each arrives \\
15 & IV & Funding Stewardship and Legislative Engagement Plan & Where the money goes, what may be lobbied for, and with whose funds \\
\end{tabularx}
\end{mvtable}
\tabcapl{tab:index}{The fifteen documents. Document 13 is the load-bearing
addition: a reader who deletes every model and every robot from the site still
holds a complete conventional Phase 1 requirements manual.}

\section{The Site at a Glance}

The parameters below are fixed once, here. Documents 05 through 14 cite this
table rather than restating a dimension, so a single number cannot appear in two
documents with two values.

\begin{mvtable}
\begin{tabularx}{\textwidth}{>{\raggedright\arraybackslash}p{4.7cm} >{\raggedright\arraybackslash}p{4.4cm} >{\raggedright\arraybackslash}X}
\headrow \thead{Parameter} & \thead{Value} & \thead{Where it is used}\\
\midrule
Jurisdiction & City of San Diego, La Jolla Community Plan area & Documents 05, 10 \\
Distance to the intended partner & About 1.6 miles & Documents 11, 13 \\
Gross floor area & 12,400 square feet & Documents 05, 08 \\
Clinical floor area & 7,850 square feet & Documents 08, 09 \\
Robotic procedure suites & 2, at 720 square feet each & Documents 08, 09, 13 \\
Procedure suite clear height & 10 feet 6 inches & Document 08 \\
Exam rooms & 4 & Documents 09, 13 \\
Infusion positions & 4 & Documents 09, 13 \\
Investigational drug pharmacy & 340 square feet & Documents 08, 13 \\
Specimen processing laboratory & 260 square feet & Documents 08, 13 \\
Machine room & 180 square feet & Documents 08, 09, 11 \\
Robot types and instances & 8 types, 14 instances & Documents 06, 09, 11 \\
Parking stalls & 46 & Document 10 \\
Hours of operation & 06:00 to 20:00 weekdays; 08:00 to 14:00 Saturday & Documents 05, 12, 14 \\
Participants & Up to 18 treated, up to 30 screened & Documents 10, 13 \\
Trial design & Open-label, single-arm, 3+3 escalation & Document 13 \\
Staff & 11 roles, 3.95 award-funded full-time equivalents & Documents 06, 14 \\
Direct award & \$700,000 per year, \$3,500,000 over five years & Documents 04, 14, 15 \\
Activation target & Week 46 from lease execution & Documents 12, 14 \\
\end{tabularx}
\end{mvtable}
\tabcapl{tab:site}{Site parameters, fixed once and cited thereafter. The
leasehold address is withheld until a lease is executed; the site is identified
by community plan area and by zone in document 05.}

\section{The Author Record and the Funding Position}

ChemicalQDevice has been a California limited liability company since October
2021. Its first years produced no application: a provisional patent filing that
did not convert, forty weeks of seminars, a six-month literature review, and two
years of notebook work whose conclusion was that the quantum-inspired approach it
had been built around could not escape exponential resource growth as the
algorithms scaled. That record is stated here because it is the reason the
current method is what it is. The company reached a working method in 2024, when
multimodal large language models became able to carry text, image, and document
context together, and the first full manuscript that followed tested exactly
that capability \cite{lmmchem2024}.

Of the fifty-six works in the author's deposited corpus, thirty-six were
generated with artificial intelligence assistance in 2026 alone. Twenty of them
form the chronology that this trial rests on, from drug identification through
protocol, investigational new drug application, and funding
\cite{pdacdeck2026}. The development practice is public: demonstrations and
sessions are run open-source and online alongside other developers, so the
method is inspectable rather than asserted.

The funding position assumed throughout this package is a federal award of
\$700,000 per year for five years, \$3,500,000 in direct cost, from at least one
federal agency. Favorable responses to project applications have been received
from federal funders, three presidential recognition letters have been issued to
the chief executive, and industry recipients of the author's communications have
been responsive. Each of those facts is stated once more in document 15, with
what it does and does not establish. The budget itself is not re-derived
anywhere in this package: it is reused from the ten-application work
\cite{tenapps}.

\section{How to Read This Package}

\begin{mvtable}
\begin{tabularx}{\textwidth}{>{\raggedright\arraybackslash}p{4.0cm} >{\raggedright\arraybackslash}p{2.6cm} >{\raggedright\arraybackslash}X}
\headrow \thead{Reader} & \thead{Documents} & \thead{Why, and in what order}\\
\midrule
Legislative staff & 01, 02, 03, 04 & Read 01 first for the definitions the other three inherit, then whichever instrument is in front of you \\
City planner & 05, then 08, 09, 10 & 05 states the entitlement; the three that follow are the conditions a permit would attach \\
State surveyor & 06, then 12, 14 & 06 is the survey instrument; 12 and 14 hold the evidence a survey asks for \\
FDA reviewer & 07, then 13 & 07 maps the site to federal requirements; 13 shows the conventional trial those requirements govern \\
Funder & Front matter, 13 §7, 14, 15 & The evidence, the roster, and the money, in that order \\
Site principal investigator & 13, then 12, 09 & The trial first, then the procedures, then the premises rules that constrain them \\
Prospective participant & 02, then 10 & The rights that attach to a participant, and how one gets to and from the site \\
\end{tabularx}
\end{mvtable}
\tabcapl{tab:audience}{Seven readers and the shortest path through the package
for each. No reader is expected to read all fifteen documents, which is why each
one repeats the definitions it depends on rather than referring outward.}
